\documentclass[a4paper,12pt]{article}
\usepackage{bbding,combelow,textcomp,amsfonts,amsthm,amsmath,amssymb,graphicx,color,hyperref,etoolbox}
\usepackage[utf8x]{inputenc}
\usepackage[lined,boxed,commentsnumbered]{algorithm2e}

\newtoggle{story}


%%%%%%%%%%%%%%%%%%%%%%%%%%%%%%%%%%%%%%%%%%
%% Customisation options --- update as necessary
%%%%%%%%%%%%%%%%%%%%%%%%%%%%%%%%%%%%%%%%%%

%% Course details: title, semester, instructors

\newcommand{\courseTitle}{Combinatorics II (Math 7702)}
\newcommand{\semester}{Spring 2026}
\newcommand{\instructorA}{Shagnik Das}
\newcommand{\instructorB}{Lu Yun-Chi}

%% Sheet number

\newcommand{\sheetNumber}{1}


%% Submission details

\newcommand{\dueDate}{15:30, Mar 24th.}
\newcommand{\lateWarning}{Late submissions will receive no credit.}


%% Display irrelevant footnotes?  \toggletrue for yes, \togglefalse for no

\toggletrue{story}

%%%%%%%%%%%%%%%%%%%%%%%%%%%%%%%%%%%%%%%%%%
%% End of customisation options
%%%%%%%%%%%%%%%%%%%%%%%%%%%%%%%%%%%%%%%%%%

\pdfpagewidth 8.5in
\pdfpageheight 11in
\topmargin -1in
\headheight 0in
\headsep 0in
\textheight 8.5in
\textwidth 6.5in
\oddsidemargin 0in
\evensidemargin 0in
\headheight 77pt
\headsep 0in
\footskip .75in

\makeatletter
\newcommand{\buildtitle}[4]{
\begin{flushleft}
{\large
#1
\hfill{}
#2
\par
#3
}
\end{flushleft}
\vskip 4pt
\begin{center}
{\large\bfseries#4\par}
\end{center}
\bigskip
}
\makeatother

\renewcommand{\thesection}{\normalsize\arabic{section}}

\renewcommand{\deg}{\textrm{deg}}
\newcommand{\eps}{\varepsilon}
\newcommand{\ceil}[1]{\left \lceil #1 \right \rceil}

\newcommand{\hint}{\begin{flushright} [Hint at \url{\hintURL}.] \end{flushright}}

\newcommand{\card}[1]{\left| #1 \right|}
\newcommand{\mb}[1]{\mathbb{#1}}
\newcommand{\mc}[1]{\mathcal{#1}}


\newcommand{\story}[1]{\iftoggle{story}{\footnote{#1}}{}}
\newcommand{\storymark}[1][42]{\iftoggle{story}{\footnotemark[#1]}{}}
\newcommand{\storytext}[2][42]{\iftoggle{story}{\footnotetext[#1]{#2}}{}}

\newcommand{\bonus}[2]{\paragraph{Bonus (#1 pt\ifstrequal{#1}{1}{}{s})} #2}


\begin{document}


\buildtitle{\courseTitle{}}{\semester}{\instructorA{} \hfill \instructorB{}}{Exercise Sheet \sheetNumber{}}

\vspace{-0.2in}

\begin{center}

{\bf Due date: \dueDate \\ \lateWarning}

\end{center}

\noindent You should try to solve all of the exercises below, and then submit two solutions to be graded --- each problem is worth $10$ points. Make sure to clearly justify your answers, defining everything precisely and stating any results you are using. You should then submit your solutions on NTU COOL, either as a typed PDF or a clear picture or scan of your handwritten answers.

\paragraph{Exercise 1}  Consider the following algorithm to find the minimum of a set of $2^k$ numbers.
\begin{algorithm}[H]
\medskip
\TitleOfAlgo{MIN}
\KwData{$A = \{a_1, \hdots, a_n\}$, $n = 2^k \ge 2$}
\KwResult{MIN($A$) = $\min \{a_1, \hdots, a_n\}$}
\medskip
\eIf{$n = 2$}
{\eIf{$a_1 < a_2$}{return $a_1$\;}{return $a_2$\;}}
{\For{$1 \le i \le n/2$}{set $y_i$ = MIN($\{a_{2i-1}, a_{2i}\}$)\;}
return MIN($\{y_1, \hdots, y_{n/2}\}$)\;}
\end{algorithm}

\begin{itemize}
	\item[(a)] Show that the MIN algorithm requires $n-1$ comparisons to find the minimum element, and that this is the best possible.
	\item[(b)] After running the MIN algorithm to find the minimal element, how many additional comparisons are required to find the second-smallest element?
	\item[(c)] Deduce a sorting algorithm that requires $\sim n \log_2 n$ comparisons to sort $n = 2^k$ elements.
\end{itemize}

\paragraph{Exercise 2}  Consider the following game.  I think of an integer $x$ between $1$ and $n$, and your job is to try and determine $x$.  You are allowed to ask questions of the form ``Is $x < y$?" or ``Is $x > y$?" for any $y$.
\begin{itemize}
	\item[(a)] Show that you can find $x$ with only $\ceil{\log_2 n}$ questions, and that this is best possible.
\end{itemize}
To make your job slightly harder, I am now allowed to lie to you at most $k$ times, for some constant $k$.
\begin{itemize}
	\item[(b)] How many questions do you now need to determine $x$?  Provide the best lower and upper bounds that you can find.
\end{itemize}

\paragraph{Exercise 3}  Given a connected graph $G$ and an arbitrary vertex $v_0 \in V(G)$, show that the breadth-first search starting at $v_0$ returns a tree $T_B$ that preserves distances\footnote{In an unweighted graph $G = (V,E)$, the \emph{distance} $d_G(u,v)$ between vertices $u, v \in V$ is the minimum length of a path from $u$ to $v$.  If $u$ and $v$ are in separate connected components, we may take their distance to be infinite.} to $v_0$; that is, for every $v \in V(G)$, $d_G(v_0, v) = d_{T_B}(v_0, v)$.

\paragraph{Solution:} 
Note that since \(G\) is connected, so every \(v \in V(G)\) has \(v \in T_B\).   
Now suppose $v \in V(G)$ satisfies $d_G(v_0,v)=d$.

Let
\[
P: v_0 \to v_1 \to \cdots \to v_d = v
\]
be the path in $T_B$ from $v_0$ to $v$. We claim that
\[
d_G(v_0,v_i)=d_{T_B}(v_0,v_i)=i \qquad \text{for all } 0\le i\le d .
\]

We prove this by induction on $i$.

\textbf{Base case:} $i=0$. This is trivial since $v_0=v_0$ and
\[
d_G(v_0,v_0)=d_{T_B}(v_0,v_0)=0 .
\]

\textbf{Inductive step:} Suppose the statement is true for all $0\le i<k$. Then
\[
v_0 \to v_1 \to \cdots \to v_{k-1}
\]
is a shortest path from $v_0$ to $v_{k-1}$. Hence
\[
v_0 \to v_1 \to \cdots \to v_{k-1} \to v_k
\]
is a shortest path from $v_0$ through \(v_{k-1}\) to $v_k$ of length $k$.

Now suppose that $d_G(v_0,v_k)<k$. Then there exists a path $P'$ from $v_0$ to $v_k$ of length $d'<k$. Let $u$ be the vertex just before $v_k$ on the path $P'$, so that
\[
{u,v_k}\in P'.
\]
Note that $u\ne v_{k-1}$, otherwise we would obtain a shorter path from $v_0$ through $v_{k-1}$ to $v_k$, which is impossible.

Since $d'<k$, by the induction hypothesis the path in $T_B$ from $v_0$ to $u$ is a shortest path, and its length is $d'-1$. Hence $u$ must be added to the set $W$ (the set of explored vertices in breadth-first search) earlier than $v_{k-1}$.

Therefore,
\[
d_{T_B}(v_0,v_k) \le d_{T_B}(v_0,u)+1=d_G(v_0,u)+1=(d'-1)+1=d'<k,
\]
which shows that
\[
v_0 \to v_1 \to \cdots \to v_{k-1} \to v_k
\]
is not the path in $T_B$ from $v_0$ to $v_k$, a contradiction.

Hence $d_G(v_0,v_k)\ge k$. Since we already have a path of length $k$ from $v_0$ to $v_k$, we conclude that
\[
d_G(v_0,v_k)=k=d_{T_B}(v_0,v_k).
\]

This completes the induction, and therefore for every vertex $v$ in the same connected component as $v_0$,
\[
d_G(v_0,v)=d_{T_B}(v_0,v).
\]



\paragraph{Exercise 4}  Given a weighted graph $G$ with $m$ edges, order the edges by weight so that $w(e_1) \ge w(e_2) \ge \hdots \ge w(e_m)$.  Set $G_0 = G$.  At step $1 \le t \le m$, if the graph $G_{t-1} - e_t$ is connected, delete $e_t$, and set $G_t = G_{t-1} - e_t$.  Otherwise keep $e_t$ and set $G_t = G_{t-1}$.

Prove that the final graph $G_m$ is a minimum weight spanning tree.

% \newpage

% \section*{Hints}  The following QR codes contain hints for some of the homework exercises.  You should be able to decode them using any QR code scanner, including this one: \url{https://zxing.org/w/decode.jspx}.

% \paragraph{Exercise ?} Also available at \url{https://i.ibb.co/???????/S??E?.png}.

\begin{center}
% \includegraphics[scale=0.5]{Hints/S??E?.png}
\end{center}

\end{document}